
@@M0@@
les @@M1@@
étant les classes de faisceaux inversibles : pour le montrer, on se ramène par le raisonnement de 3.1 au cas d'un fibré projectif sur @@M2@@. Or ce dernier cas est conséquence immédiate de 1.13. L'élément @@M3@@ considéré est dans @@M4@@ d'après ce qui précède, et par conséquent @@M5@@. Or @@M6@@; donc d'après 6.6, @@M7@@. Comme @@M8@@ fait partie d'une base de @@M9@@ sur @@M10@@, ceci entraine @@M11@@. Comme @@M12@@ est injectif (3.1), on a finalement @@M13@@.

\textbf{Problème 6.10.} Sous les hypothèses de 6.9, est-il vrai que pour tout @@M14@@ on ait : @@M15@@ ?

\subsection*{BIBLIOGRAPHIE}
\begin{enumerate}[label={\textup{[\arabic*]}}]
\item Bass H., K-Theory and stable algebra, Publ. Math. de l'IHES, n°22.
\item Chevalley C., Séminaire 1958 : Anneaux de Chow et applications.
\item Hakim M., Topos localement annelés et schémas relatifs, thèse (en préparation).
\item Hartshorne R., Residues and duality, Lecture Notes in Mathematics n°20, Springer-Verlag.
\item Serre J.-P., Le groupe de Grothendieck des schémas en groupes réductifs, Publications Mathématiques de l'IHES, n°34.
\end{enumerate}



% Expose VII additions
\providecommand{\Kosz}{\mathrm K}
\providecommand{\Proj}{\operatorname{Proj}}
\providecommand{\Sym}{\operatorname{Sym}}
\providecommand{\gr}{\operatorname{gr}}
\providecommand{\Tor}{\operatorname{Tor}}
\providecommand{\Spec}{\operatorname{Spec}}
\providecommand{\Imf}{\operatorname{Im}}
\providecommand{\Ker}{\operatorname{Ker}}
\providecommand{\Id}{\operatorname{Id}}
\providecommand{\rad}{\operatorname{rad}}
\providecommand{\cO}{\mathcal O}
\providecommand{\cS}{\mathcal S}
\providecommand{\cJ}{\mathcal J}
\providecommand{\cN}{\mathcal N}
\providecommand{\bfP}{\mathbf P}
\providecommand{\bfA}{\mathbf A}


